\section{Perpendicular standing spin waves in AuPt/GdFeCo induced by spin-orbit torques - \href{https://doi.org/10.1038/s42005-025-02433-2}{Comm. Phys.,8, 516 (2025)-Nature}}

\section*{Main points from the Paper}
\subsection*{\underline{I. Introduction:}}
    \begin{enumerate}[1.]
        \item GdFeCo has a ferrimagnetic compositions, their spin dynamics in the present regime are effectively ferromagnetic (FM).
        \item They have observed the SOT, generated from the adjacent AuPt layer induce perpendicular standing spin waves in GdFeCo, with resonant modes observed up to the third order. 
        \item Uncapped films exhibit oxidation-induced degradation that suppresses spin-wave confinement and to avoid that they have used Ta as the capping layer. 
        \item Ferrimagnetic materials (FiMs) can also be used to provide most of the things which are useful in device chips which require a faster switching speeds and reduced saturation magnetization [\textcolor{red}{first lines are good}].
        \item The combined properties of FM and AFM are found in the FiMs: low net saturation magnetization, fast magnetic dynamics driven by antiferromagnetic precession, and the ease of magnetization manipulation and detection typically seen in ferromagnets.
        \item Because of the presence of the forced precession of antiparallel spins in FiMs, there is formation of perpendicular standing spin-waves(PSSW) modes, where spins oscillate perpendicular to the film surface.
        \item The formation of PSSW modes was observed recently in AFM or FiM systems in connection with FM which is attributed to the interlayer exchange coupling with the adjacent FM layer. 
        \item ST-FMR is used to observe the PSSW modes in AuPt/GdFeCo.
        \item In the case of uniform precession of the magnetic moments $(p=0)$ or PSSW modes when spin waves witht non-zero wave vectors satisfy the confinement condition \boxed{$(k=p\pi/d)$} where $d$ is the film thickness.
        \item The Ta capping layer was used as capping layer to reduce the oxidation and enabled us to stabilize the PSSW modes up to the $p=3$. And the samples without the Ta capping, the ST-FMR spectra with PSSW modes were limited to the $p=2$ and a reduced saturation magnetization compared to those with the Ta capping layer.
        \item These findings provide valuable insights for the development of FiM-based SOT-MRAM and logic devices, which benefit from the unique properties of FiMs.
    \end{enumerate}
\subsection*{\underline{II. Results and Discussion:}}
\begin{enumerate}[1.]
    \item VSM is used to measure the fundamental magnetic properties of the system, and saturation magnetization $(M_s)$ and effective anisotropy energy density $(K_{eff})$ as a function of GdFeCo thickness is calculated. \footnote{See the supplementary note 1 of the paper for detailed calculations of these.}
    \item Surface effects was taken as the reason why there is larger saturation magnetization $(M_s)$. On increasing thickness, the bulk properties becomes more dominant enhancing ferrimagnetic ordering and a reduction of overall magnetization.
    \item Suppression of oxidation significantly affects the magnetization, particularly at lower thicknesses.[\textcolor{red}{Important Line}]
    \item Surface anisotropy weakens as $t_{GdFeCo}$ increases, leading to bulk anisotropy becoming dominant contribution. It is consisted with $M_s$ results.
    \item No PMA has been observed across the entire thickness range.
    \item The spin currents generated by the Spin Hall Effect (SHE) from the strong SOC in the adjacent AuPt layer and the induced $H_{Oe}$ produced by the RF current drives the magnetization dynamics in the GdFeCo layer, where the antiparallel alignment of magnetic moments results in varying precession amplitudes and directions among individual atoms.
    \item The excited spin-wave modes generated by this dynamic behavior exhibit the PSSW modes when teh wave vectors $(k)$ satisfy the confinement condition, $k=p\pi/t_{GdFeco}$.
    \item The Ta layer prevents the formation of FeO, CoO, or GdO hence the 3rd mode $(p=3)$ is only present in the case of Ta capping.
    \item Experimental evidence of presence of 3rd PSSW mode in the presence of the Ta layer was observed. This enhances the surface pinning and spin-wave reflection, promoting the formation of higher-order PSSW modes.
    \item The sharp increase in $(A_{ex})$ at $t_{GdFeCo}=80$ nm is attributed to the dominant bulk properties in the GdFeCo layer. This irregular behavior suggests that spin-wave propagation and PSSW formation are hindered by spin-wave absorption and reduced spin-wave reflection due to surface oxidation.
    \item The interfacial engineering after adding Ta as the capping layer, stabilizes the ferrimagnetic dynamics, promoting the formation of higher-order spin-wave modes.
    \item The dynamic pinning model and introduced the wave vector$k_p$ defined as $k_p=\left[\frac{(p+\delta)\pi}{t_{GdFeCo}}=\frac{2\pi}{\lambda_p}\right]$, where $\lambda_p$ is the wavelength in the GdFeCo layer, accounting for each mode  $p$ and the pinning parameter $\delta$ was used by the authors to quantify the degree of pinning at the surface. [\textcolor{red}{ref. 49-51, from this paper}].
    \item When the Ta capping was done on the system here, this pinning model specified that there is strong interfacial confinement and in the case of not Ta capping, there was significantly lower confinement.
    \item Conclusively, the Ta capping enhances the interfacial confinement, supporting the formation of stable and higher-order PSSW modes.
    \item When magnetic field polarity measurements were carried out, the 1st PSSW mode was only visible under a positive external magnetic field. In contrast, the Ta-capped samples exhibit symmetrical resonance under both positive and negative external fields. 
    \item The oxidation of the Gd leads to the asymmetrically aligned local magnetic moments, breaking the mirror symmetry with respect to the external magnetic field.
    \item Potential mechanisms may include non-reciprocal spin-wave propagation in magnetization-graded systems or the influence of the Dzyaloshinskii-Moriya interaction (DMI) [\textcolor{red}{ref. 56-59 from this paper}]. 
\end{enumerate}
\subsection*{\underline{III. Summary:}}
\begin{enumerate}
    \item An investigation of the formation of PSSW modes in ferrimagnetic GdFeCo layers drive by SOTs generated by the adjacent high SOC material, AuPt was carried out.
    \item The ST-FMR spectra explains the induced antiparallel magnetization precession after the contact of AuPt which injected the SOT.
    \item Oxidation at the surface of GdFeCo was shown to significantly affect the key properties such as $M_s,A_{ex}$, and overall quality of deposited GdFeCo layer.
    \item The pinning effects at the interfaces were evaluated using a dynamic pinning model, demonstrating that the Ta-capped samples provided strong interfacial confinement for PSSW modes, while the uncapped samples exhibited relaxed boundary conditions and weaker pinning.
\end{enumerate}
